\chapter{Period Problems}
\index{Period Problems}

\medskip
\noindent\textit{This chapter is for general education. People who menstruate can seek confidential assessment for symptoms that affect daily life.}

\section*{Common problems}
Period problems include painful periods, heavy or prolonged bleeding, irregular or missed periods, premenstrual symptoms, menstrual migraine, and bleeding after menopause. Patterns vary with puberty, pregnancy, contraception, perimenopause, stress, weight change, exercise, medicines, and health conditions.

\section*{Possible causes}
Causes include normal hormonal variation, pregnancy, polycystic ovary syndrome, thyroid disease, fibroids, endometriosis, adenomyosis, pelvic infection, bleeding disorders, and medicines or contraceptive devices. A change from a person’s usual pattern deserves attention, especially with pain or bleeding between periods.

\section*{Assessment and treatment}
Keep a diary of dates, bleeding, pain, medicines, and other symptoms. Assessment may include a pregnancy test, examination, blood tests, and ultrasound or other tests. Treatment depends on the cause and may include heat, suitable pain relief, hormonal treatment, medicines to reduce bleeding, or specialist care. Check pain medicines with a pharmacist if you have asthma, stomach, heart, kidney, or liver problems.

\section*{When to Seek Medical Care}
Seek urgent care for severe or unusual pelvic pain, fainting, possible pregnancy with pain or bleeding, or bleeding that rapidly soaks through products and causes weakness or dizziness. Arrange review for periods that are very painful, heavy, irregular, last longer than usual, stop unexpectedly, cause anaemia symptoms, or include bleeding after menopause.

\section*{Sources and further reading}
\begin{itemize}
\item NHS. \textit{Period problems}. \url{https://www.nhs.uk/conditions/periods/period-problems/}
\item NHS. \textit{Period pain}. \url{https://www.nhs.uk/symptoms/period-pain/}
\end{itemize}
